\documentclass[11pt,a4paper]{article}

% ---------- Packages ----------
\usepackage[T1]{fontenc}
\usepackage[utf8]{inputenc}
\usepackage{lmodern}
\usepackage[margin=25mm]{geometry}
\usepackage{setspace}
\usepackage{amsmath,amssymb,amsfonts,amsthm,mathtools}
\usepackage{mathrsfs}
\usepackage{bm}
\usepackage{enumerate}
\usepackage{hyperref}
\usepackage{cleveref}
\usepackage{cite}

% ---------- Line spacing ----------
\onehalfspacing

% ---------- Hyperref ----------
\hypersetup{
  colorlinks=true,
  linkcolor=blue,
  citecolor=blue,
  urlcolor=blue,
  pdftitle={Title of the Paper},
  pdfauthor={Author Name}
}

% ---------- Theorem environments ----------
\theoremstyle{plain}
\newtheorem{theorem}{Theorem}[section]
\newtheorem{proposition}[theorem]{Proposition}
\newtheorem{lemma}[theorem]{Lemma}
\newtheorem{corollary}[theorem]{Corollary}

\theoremstyle{definition}
\newtheorem{definition}[theorem]{Definition}
\newtheorem{assumption}[theorem]{Assumption}
\newtheorem{example}[theorem]{Example}
\newtheorem{remark}[theorem]{Remark}

% ---------- Custom commands ----------
\newcommand{\R}{\mathbb{R}}
\newcommand{\Q}{\mathbb{Q}}
\newcommand{\Z}{\mathbb{Z}}
\newcommand{\N}{\mathbb{N}}
\newcommand{\C}{\mathbb{C}}
\newcommand{\F}{\mathbb{F}}

\newcommand{\eps}{\varepsilon}
\newcommand{\To}{\Rightarrow}
\newcommand{\into}{\hookrightarrow}
\newcommand{\onto}{\twoheadrightarrow}

\newcommand{\GL}{\mathrm{GL}}
\newcommand{\SL}{\mathrm{SL}}
\newcommand{\PSL}{\mathrm{PSL}}
\newcommand{\Aut}{\mathrm{Aut}}
\newcommand{\Hom}{\mathrm{Hom}}
\newcommand{\End}{\mathrm{End}}
\newcommand{\Ker}{\mathrm{Ker}}
\newcommand{\Img}{\mathrm{Im}}
\newcommand{\tr}{\mathrm{tr}}
\newcommand{\rank}{\mathrm{rank}}
\newcommand{\ord}{\mathrm{ord}}

\newcommand{\abs}[1]{\left\lvert #1 \right\rvert}
\newcommand{\norm}[1]{\left\lVert #1 \right\rVert}
\newcommand{\set}[1]{\left\{ #1 \right\}}
\newcommand{\angles}[1]{\left\langle #1 \right\rangle}
\newcommand{\paren}[1]{\left( #1 \right)}
\newcommand{\bracket}[1]{\left[ #1 \right]}

% ---------- Title ----------
\title{\textbf{Title of the Paper}}
\author{
  Author Name\\
  \small Department of Mathematics, University Name\\
  \small \texttt{author@example.com}
}
\date{\today}

% ---------- Document ----------
\begin{document}

\maketitle

\begin{abstract}
We investigate a mathematical problem related to [insert topic].
The main purpose of this paper is to establish [insert main result].
Our approach is based on [insert method], and several examples are provided
to illustrate the theory.
\end{abstract}

\section{Introduction}

The purpose of this paper is to study [insert topic].
This subject has been investigated by many authors in various contexts.
Our main result can be regarded as a refinement of earlier work on
[insert related topic].

The paper is organized as follows.
In \Cref{sec:preliminaries}, we collect basic definitions and notation.
In \Cref{sec:main-results}, we state and prove the main theorem.
Finally, in \Cref{sec:examples}, we discuss several examples and remarks.

\section{Preliminaries}\label{sec:preliminaries}

We begin with the following definition.

\begin{definition}
Let $X$ be a set and let $f : X \to X$ be a map.
We say that $f$ has property \emph{(P)} if
\[
  f(x) = x
\]
for every $x \in X$ satisfying [insert condition].
\end{definition}

The following lemma will be used repeatedly.

\begin{lemma}\label{lem:basic}
Let $V$ be a vector space over $\R$, and let $T \in \End(V)$.
If $T$ is diagonalizable, then
\[
  V = \bigoplus_{\lambda} E_{\lambda},
\]
where $E_{\lambda}$ denotes the eigenspace corresponding to the eigenvalue $\lambda$.
\end{lemma}

\begin{proof}
This is a standard consequence of the definition of diagonalizability.
\end{proof}

\section{Main Results}\label{sec:main-results}

We now state the principal theorem of this paper.

\begin{theorem}\label{thm:main}
Let $G$ be a group acting on a set $X$.
Assume that [insert assumptions].
Then the following assertions hold:
\begin{enumerate}[(i)]
  \item the action of $G$ on $X$ satisfies [insert property];
  \item for every $x \in X$, the stabilizer subgroup $G_x$ has [insert property];
  \item there exists a natural bijection between [insert objects].
\end{enumerate}
\end{theorem}

\begin{proof}
We divide the proof into several steps.

\medskip
\noindent
\textbf{Step 1.}
First, we prove that [insert claim].
By \Cref{lem:basic}, we have
\[
  V = \bigoplus_{\lambda} E_{\lambda}.
\]

\medskip
\noindent
\textbf{Step 2.}
Next, we show that [insert claim].
For any $x \in X$, one obtains
\[
  \phi(x) = \sum_{n=0}^{\infty} a_n x^n,
\]
whenever the series converges.

\medskip
\noindent
\textbf{Step 3.}
Combining the previous steps, we conclude that the desired statement follows.
\end{proof}

As an immediate consequence, we obtain the following.

\begin{corollary}
Under the assumptions of \Cref{thm:main}, the object [insert object]
is uniquely determined up to isomorphism.
\end{corollary}

\begin{proof}
This follows directly from \Cref{thm:main}.
\end{proof}

\section{Examples and Remarks}\label{sec:examples}

We present a simple example.

\begin{example}
Let $G = \SL_2(\Z)$.
Then $G$ acts naturally on the upper half-plane by fractional linear transformations:
\[
  z \mapsto \frac{az+b}{cz+d}
  \qquad
  \left(
    \begin{pmatrix}
      a & b\\
      c & d
    \end{pmatrix}
    \in \SL_2(\Z)
  \right).
\]
This provides a classical source of examples in the theory of modular forms.
\end{example}

\begin{remark}
The argument given above can be adapted to more general settings,
for example to [insert generalization].
\end{remark}

\section*{Acknowledgments}

The author would like to thank [insert names] for valuable discussions.

\begin{thebibliography}{99}

\bibitem{Lang}
S.~Lang,
\textit{Algebra},
Graduate Texts in Mathematics, Springer, 2002.

\bibitem{Serre}
J.-P.~Serre,
\textit{A Course in Arithmetic},
Graduate Texts in Mathematics, Springer, 1973.

\bibitem{Hartshorne}
R.~Hartshorne,
\textit{Algebraic Geometry},
Graduate Texts in Mathematics, Springer, 1977.

\end{thebibliography}

\end{document}